\subsection{Hauptprogramm des IPv4-Clients}

\begin{frame}[fragile]{Hauptprogramm des IPv4-Clients}
  \begin{cppcode}[autogobble=true]
    #include "EchoClientContextFactory.h"
    #include <core/SNodeC.h>
    #include <iostream>
    #include <net/in/stream/legacy/SocketClient.h>
    #include <string.h>
    #include <string>

    int main(int argc, char* argv[]) {
      core::SNodeC::init(argc, argv);

      using EchoClient = net::in::stream::legacy::SocketClient<EchoClientContextFactory>;
      using SocketAddress = EchoClient::SocketAddress;

      EchoClient echoClient;

      echoClient.connect("localhost", 8001, [](const SocketAddress& socketAddress, int err) -> void {
        if (err == 0) {
          std::cout << "Success: Echo connected to " << socketAddress.toString() << std::endl;
        } else {
          std::cout << "Error: Echo client connected to " << socketAddress.toString()
          << ": " << strerror(err) << std::endl;
        }
      });

      return core::SNodeC::start();
    }
  \end{cppcode}
\end{frame}
